\section{Additional implementation and evaluation details}
\label{app:settings}
\subsection{Monitoring and geometry}
The LJ fit uses column-normalized least squares on a $3\times2$
matrix. It requires rank two, a normalized condition number no greater
than $10^4$, positive fitted coefficients, and a fitted
$\sigma=(\alpha/\beta)^{1/6}$ between 0.2 and 10~\AA.
The residual denominator uses $\delta=10^{-8}$ in the force units
of Eq.~\eqref{eq:lj}; invalid fits receive residual one.
The check concerns compatibility of one force vector with a local
pair-force basis. It does not test energy conservation, a full force
Jacobian, or a transferable LJ parameterization across atoms.

The diversity fingerprint comprises 16-bin geometric summaries of
pair-distance density, atomic-number-weighted density, and weighted
bin means, giving 48 components with an 8~\AA\ cutoff.
The selector compares fingerprint distances against the configured
0.1 threshold. In a pending event, later triggers merge hotspot and
reason records without repeatedly extending the original post-event
window. Both completed and truncated events preserve completeness
information. These details describe the implementation; the pilot
does not separately estimate each design choice's causal contribution.

\subsection{Training and evaluation conventions}
The static study uses a minimum learning rate of $2\times10^{-5}$,
50 warm-up optimizer steps, one attention head, maximum angular momentum one,
and an average-neighbor setting of 12. Closed-loop fine-tuning uses
five warm-up optimizer steps. The bundled configuration snapshots contain
the remaining recorded settings and their inheritance references.
They are provenance records, not standalone configurations for
launching new calculations from this manuscript directory.

Reported model force errors pool Cartesian components across the
evaluated atoms. Fragment maximum errors instead use the Euclidean
norm of each atom's three-component error. Local-label force weights
apply during training, not as a way to hide atoms in the full-parent
evaluation metric. The atomwise fidelity data retain repeated
atom--fragment observations, since a parent atom can appear in
more than one independently extracted site.

For the closed-loop trajectory, all 401 saved frames have nearest-image
separation above 20.7~\AA. Across the initial/final frozen-model runs,
the 1206 checked frames have minimum separation approximately
18.76~\AA. None violates the nominal 16~\AA\ force-dependence bound,
although 82 frames fall within a further 4~\AA\ caution margin.
These checks assess the recorded finite trajectories only.

\subsection{Exploratory protocol history}
Earlier static-pool configurations did not exercise the intended
PHAL acquisition path. The archived v2 run records zero PHAL-local
and PHAL-full acquisitions because the light-stage trigger does not
activate the higher stages. Its zero-acquisition ablations therefore
do not measure the benefit of local learning. The executed v2c
comparison follows corrections to static-frame event handling and
candidate mapping and explicitly evaluates every stage. We report
that configuration as an exploratory comparison, not as a
prospectively fixed experiment. Figures contain directly evaluated
rounds; interpolated budget metrics in the archived reports are not
presented as observations.

\section{Artifact contents and reproduction}
The source package supplies six figures in vector PDF and SVG formats
and PNG previews. Their inputs are frozen in
\texttt{data/plot\_data.json}; source records, configuration snapshots,
and hashes accompany them. The plotting script regenerates the figures
and two numerical tables without invoking GMD, CP2K, or GMD\_SGT.
A separate bundle validator checks file references, citation keys,
snapshot hashes, and numerical consistency against the per-atom
and per-run records. These checks do not replace author review of
the scientific interpretation or a final typesetting inspection.
